\documentclass[a4j]{jarticle}
\usepackage{ascmac}
\usepackage[dvipdfmx]{graphicx}
\usepackage{listings,jvlisting, jlisting}

\lstset{
  basicstyle={\ttfamily},
  identifierstyle={\small},
  commentstyle={\smallitshape},
  keywordstyle={\small\bfseries},
  ndkeywordstyle={\small},
  stringstyle={\small\ttfamily},
  frame={tb},
  breaklines=true,
  columns=[l]{fullflexible},
  numbers=left,
  xrightmargin=0zw,
  xleftmargin=3zw,
  numberstyle={\scriptsize},
  stepnumber=1,
  numbersep=1zw,
  lineskip=-0.5ex
}

\begin{document}

\vfill
\begin{screen}
  \begin{center}
    \huge\bf 実験　第1回レポート\\
    \vspace{5mm}
    \large\bf
    学籍番号　2022311042 \quad 名前 神野友哉\\
    \today
  \end{center}
\end{screen}
\thispagestyle{empty}


\section{実験1-1}
実験1-1について画像フォーマットと画像データの格納方法の関係について詳しく考察せよ.
\subsection{考察}


\newpage
\begin{figure}[htbp]
  \begin{center}
    \includegraphics[width=70mm]{image01.png}
  \end{center}
  \caption{image01.pgm}
  \label{fig:01}
\end{figure}
\begin{figure}[htbp]
  \begin{minipage}{0.5\hsize}
    \begin{center}
      \includegraphics[width=70mm]{image02.png}
    \end{center}
    \caption{image02.pbm}
    \label{fig:02}
  \end{minipage}
  \begin{minipage}{0.5\hsize}
    \begin{center}
      \includegraphics[width=70mm]{image03.png}
    \end{center}
    \caption{image03.jpg}
    \label{fig:03}
  \end{minipage}
\end{figure}

\newpage
\section{実験1-2}
実験1-2についてp=0.3とp=0.7の出力画像を示し、p-タイル法の特徴を詳しく考察せよ.
\subsection{考察}



\newpage
\begin{figure}[htbp]
  \begin{center}
    \includegraphics[width=70mm]{image1-2-1.png}
  \end{center}
  \caption{image1-2-1.pgm}
  \label{fig:121}
\end{figure}
\begin{figure}[htbp]
  \begin{minipage}{0.5\hsize}
    \begin{center}
      \includegraphics[width=70mm]{image1-2-1-30out.png}
    \end{center}
    \caption{黒 0.3}
    \label{fig:12130}
  \end{minipage}
  \begin{minipage}{0.5\hsize}
    \begin{center}
      \includegraphics[width=70mm]{image1-2-1-70out.png}
    \end{center}
    \caption{黒 0.7}
    \label{fig:12170}
  \end{minipage}
\end{figure}

\newpage
\begin{figure}[htbp]
  \begin{center}
    \includegraphics[width=70mm]{image1-2-2.png}
  \end{center}
  \caption{image1-2-2.pgm}
  \label{fig:122}
\end{figure}
\begin{figure}[htbp]
  \begin{minipage}{0.5\hsize}
    \begin{center}
      \includegraphics[width=70mm]{image1-2-2-30out.png}
    \end{center}
    \caption{黒 0.3}
    \label{fig:12230}
  \end{minipage}
  \begin{minipage}{0.5\hsize}
    \begin{center}
      \includegraphics[width=70mm]{image1-2-2-70out.png}
    \end{center}
    \caption{黒 0.7}
    \label{fig:12270}
  \end{minipage}
\end{figure}

\newpage
\begin{figure}[htbp]
  \begin{center}
    \includegraphics[width=70mm]{image1-2-3.png}
  \end{center}
  \caption{image1-2-3.pgm}
  \label{fig:123}
\end{figure}
\begin{figure}[htbp]
  \begin{minipage}{0.5\hsize}
    \begin{center}
      \includegraphics[width=70mm]{image1-2-3-30out.png}
    \end{center}
    \caption{黒 0.3}
    \label{fig:12330}
  \end{minipage}
  \begin{minipage}{0.5\hsize}
    \begin{center}
      \includegraphics[width=70mm]{image1-2-3-70out.png}
    \end{center}
    \caption{黒 0.7}
    \label{fig:12370}
  \end{minipage}
\end{figure}


\newpage
\begin{lstlisting}[caption= 実験1-2,label=c-1-2]

\end{lstlisting}
\lstinputlisting[language={Tex}]{question4next}

\end{document}
